\clearpage
\VerdierPhysicalPageAnchor{0259}%
\begingroup
\small
\setlength{\parindent}{0pt}%
\setlength{\parskip}{0pt}%
\setlength{\emergencystretch}{2em}%
\newcommand{\VerdierTermRow}[3]{%
  \VerdierSemanticAnchor{#1}%
  \par\noindent\hangindent=1.5em\hangafter=1
  #2。\enspace\textup{#3。}\par\smallskip}

\VerdierTermRow{idxt-p259-r01}{$M$ 型複体}{第 III 章、1.1.4}
\VerdierTermRow{idxt-p259-r02}{入射型（それぞれ射影型）の複体}{第 III 章、1.1.5}
\VerdierTermRow{idxt-p259-r03}{加法圏の複体}{第 I 章、2.1.5}
\VerdierTermRow{idxt-p259-r04}{Cartan--Eilenberg 入射複体}{第 III 章、4.6.1}
\VerdierTermRow{idxt-p259-r05}{加法圏の対象からなる $n$-重複複体}{第 I 章、2.1.2}
\VerdierTermRow{idxt-p259-r06}{加法圏の $n$-重複複体}{第 I 章、2.1.3}
\VerdierTermRow{idxt-p259-r07}{$n$-重複複体（$\varphi$-総和可能）}{第 I 章、2.1.7}
\VerdierTermRow{idxt-p259-r08}{総和可能な $n$-重複複体}{第 I 章、2.1.7}
\VerdierTermRow{idxt-p259-r09}{$n$-重複複体に付随する単純複体}{第 I 章、2.2.4}
\VerdierTermRow{idxt-p259-r10}{加法圏の単純複体}{第 I 章、2.1.5}
\VerdierTermRow{idxt-p259-r11}{ホモトピックな複体}{第 I 章、2.5.11}
\VerdierTermRow{idxt-p259-r12}{次数 $\sigma$ の成分（$n$-重複複体）}{第 I 章、2.1.3}
\VerdierTermRow{idxt-p259-r13}{次数 $\sigma$ の成分（$n$-重複複体の射）}{第 I 章、2.1.5}
\VerdierTermRow{idxt-p259-r14}{$\beta$-射の合成（捩れ $G$-関手）}{第 I 章、1.4.7}
\VerdierTermRow{idxt-p259-r15}{複体の射の錐}{第 I 章、3.1.2}
\VerdierTermRow{idxt-p259-r16}{$\delta$-関手}{第 III 章、1.3.1}
\VerdierTermRow{idxt-p259-r17}{標準 $\delta$-関手}{第 III 章、1.3.2}
\VerdierTermRow{idxt-p259-r18}{射の分解}{第 II 章、1.2.8}
\VerdierTermRow{idxt-p259-r19}{射の次数}{第 I 章、1.1.1}
\VerdierTermRow{idxt-p259-r20}{複体の第二標準フィルトレーション}{第 III 章、4.3.1}
\VerdierTermRow{idxt-p259-r21}{複体に付随する $D(\mathcal A)$ に値をとる第二スペクトル対象}{第 III 章、4.3.1}
\VerdierTermRow{idxt-p259-r22}{第二標準スペクトル対象}{第 III 章、4.3.1、4.3.4}
\VerdierTermRow{idxt-p259-r23}{複体に関する関手の第二スペクトル系列}{第 III 章、4.4.5}
\VerdierTermRow{idxt-p259-r24}{$\varepsilon$-可換図式}{第 I 章、1.4.3}
\VerdierTermRow{idxt-p259-r25}{次数 $\sigma$ の微分、方向 $i$（$n$-重複複体）}{第 I 章、2.1.3}
\VerdierTermRow{idxt-p259-r26}{スペクトル対象に関する微分}{第 II 章、4.2.5}
\endgroup
